\documentclass[12pt, a4paper]{article}
\usepackage[UTF8]{ctex}

\usepackage{amsmath, amssymb, amsthm}
\usepackage{geometry}
\usepackage{fancyhdr}
\usepackage{enumitem}
\usepackage{xcolor}
\usepackage{hyperref}

\geometry{margin=1in}
\pagestyle{fancy}
\fancyhf{}
\rhead{常微分方程}
\lhead{李生丹}
\rfoot{第 \thepage 页}

\newcommand{\problem}[1]{\subsection*{题 #1}}
\newcommand{\solution}{\paragraph{解.}}

\definecolor{solutioncolor}{RGB}{0, 100, 0}

\begin{document}

\begin{center}
    {\LARGE\bfseries 常微分方程 第 2 次作业} \\[0.5em]
    {\large 常微分方程} \\[0.3em]
    李生丹 \quad | \quad 2026-10-05
\end{center}

\hrule
\vspace{1em}
\problem{1}

求解初值问题：
$$y' + 2y = e^{-x}, \quad y(0) = 1$$

\solution

解析得到微分方程: $2 y{\left(x \right)} + \frac{d}{d x} y{\left(x \right)} - e^{- x} = 0$

SymPy dsolve 求解（1 阶）

代入初始条件 $y{\left(0 \right)} = 1$

特解: $y = e^{- x}$

\textbf{Answer:}
\[\boxed{y = e^{- x}}\]

\vspace{1em}
\hrule
\vspace{1em}

\problem{2}

求二阶常系数齐次方程的通解：
$$y'' - 5y' + 6y = 0$$

\solution

解析得到微分方程: $6 y{\left(x \right)} - 5 \frac{d}{d x} y{\left(x \right)} + \frac{d^{2}}{d x^{2}} y{\left(x \right)} = 0$

SymPy dsolve 求解（2 阶）

通解: $y = \left(C_{1} + C_{2} e^{x}\right) e^{2 x}$（含任意常数 C1）

\textbf{Answer:}
\[\boxed{y = \left(C_{1} + C_{2} e^{x}\right) e^{2 x}}\]

\vspace{1em}
\hrule
\vspace{1em}

\problem{3}

求解初值问题：
$$y'' + 4y' + 4y = 0, \quad y(0) = 1, \quad y'(0) = 0$$

\solution

解析得到微分方程: $4 y{\left(x \right)} + 4 \frac{d}{d x} y{\left(x \right)} + \frac{d^{2}}{d x^{2}} y{\left(x \right)} = 0$

SymPy dsolve 求解（2 阶）

代入初始条件 $y{\left(0 \right)} = 1$, $\left. \frac{d}{d x} y{\left(x \right)} \right|_{\substack{ x=0 }} = 0$

特解: $y = \left(2 x + 1\right) e^{- 2 x}$

\textbf{Answer:}
\[\boxed{y = \left(2 x + 1\right) e^{- 2 x}}\]

\vspace{1em}
\hrule
\vspace{1em}

\problem{4}

（证明题）证明：若 $y_1$ 与 $y_2$ 是二阶齐次线性常系数方程 $y'' + py' + qy = 0$ 的两个解，则它们的任意线性组合 $y = C_1 y_1 + C_2 y_2$（$C_1, C_2$ 为任意常数）也是该方程的解，并说明该结论在何种条件下可以保证 $y$ 是通解。

\solution

设线性算子 $L[y] = y'' + py' + qy$，则方程可写为 $L[y]=0$。

算子 $L$ 是线性的：对任意常数 $C_1,C_2$ 与函数 $y_1,y_2$，有 $L[C_1 y_1 + C_2 y_2] = C_1 L[y_1] + C_2 L[y_2]$。

因为 $y_1,y_2$ 是解，所以 $L[y_1]=0,\ L[y_2]=0$。

于是 $L[C_1 y_1 + C_2 y_2] = C_1\cdot 0 + C_2\cdot 0 = 0$，即 $y=C_1 y_1 + C_2 y_2$ 也是方程的解。

要保证 $y=C_1 y_1 + C_2 y_2$ 是通解，需要 $y_1,y_2$ 线性无关，即它们的朗斯基行列式 $W(y_1,y_2)=y_1 y_2' - y_1' y_2 \neq 0$。

对于二阶齐次线性方程，两个解线性无关当且仅当 $W \neq 0$；此时 $\{y_1,y_2\}$ 构成基本解组，$y=C_1 y_1 + C_2 y_2$ 包含全部解，即为通解。

\textbf{Answer:}
\[\boxed{y = C_1 y_1 + C_2 y_2 \text{ 是解；当 } W(y_1,y_2)=y_1 y_2' - y_1' y_2 \neq 0 \text{ 时，它是通解。}}\]

\vspace{1em}
\hrule
\vspace{1em}


\vspace{2em}
\noindent\textit{Auto-generated: 4/4 problems solved by SymPy.}
\end{document}
